\documentclass[11pt,a4paper]{article}
\usepackage{fontspec}
\usepackage{xeCJK}
\setCJKmainfont{Hiragino Mincho ProN}
\setCJKsansfont{Hiragino Kaku Gothic ProN}
\usepackage[margin=2.5cm]{geometry}
\usepackage{amsmath,amssymb,amsfonts}
\usepackage{hyperref}
\usepackage{booktabs}
\usepackage{amsthm}

\newtheorem{theorem}{定理}
\newtheorem{observation}[theorem]{観察}
\newtheorem{conjecture}[theorem]{予想}

\newcommand{\Spec}{\mathrm{Spec}}

\begin{document}

\title{微細構造定数の算術的表現：\\ベルヌーイ数・素数ギャップ・ゼータ値による公式}
\author{Wright Brothers\\{\small 独立研究}}
\date{2026年4月}
\maketitle

\begin{abstract}
逆微細構造定数$\alpha^{-1} = 137.035999\ldots$に対し、ベルヌーイ数、素数ギャップ、時空次元$d$、およびリーマンゼータ関数の正整数値のみを含む閉じた表現を報告する：
$$\alpha^{-1} = \frac{2}{|B_2|} + \frac{4}{|B_4|} + g\!\left(\frac{2}{|B_2|} + \frac{4}{|B_4|}\right) + d\bigl(\zeta(2d{-}2) - \zeta(2d{-}1)\bigr).$$
ここで$g(n)$は素数ギャップ関数（$n+k$が素数となる最小の$k > 0$）、$d = 4$は時空次元。自由パラメータなしで$137.03598$を与え、実験値との差は$2.4 \times 10^{-5}$（相対誤差$1.7 \times 10^{-7}$）。公式は (i)~$12 = 2/|B_2|$（1D カシミールエネルギー分母）、(ii)~$120 = 4/|B_4|$（3D カシミールエネルギー分母）、(iii)~$5 = g(132)$（素数ギャップ）、(iv)~$4(\zeta(6)-\zeta(7))$（高次元真空揺らぎ補正）に分解される。第一原理からの導出ではなく、物理的真空が$\Spec(\mathbb{Z})$の算術構造を持つという仮説と整合する\emph{構造的観察}として提示する。公式から独立な3つの予測を導き、数値一致が偶然である確率を評価する。
\end{abstract}

\section{序論}

微細構造定数$\alpha \approx 1/137.036$は電磁相互作用の強さを支配する。物理学で最も精密に測定された無次元定数の一つであるが~\cite{Mohr2016}、その値を第一原理から予測する理論は受け入れられていない。

本研究は特定の理論的枠組み——量子真空が$\Spec(\mathbb{Z})$の算術構造を持つという仮説~\cite{paper_A,paper_F}——から問題にアプローチする。この枠組みの中で、$\alpha^{-1}$は$\Spec(\mathbb{Z})$の標準的算術不変量で表現可能か、と問う。

\section{公式}

\begin{observation}[算術的$\alpha$]
\label{obs:alpha}
\begin{equation}
  \alpha^{-1} = \frac{2}{|B_2|} + \frac{4}{|B_4|} + g\!\left(\frac{2}{|B_2|} + \frac{4}{|B_4|}\right) + d\bigl(\zeta(2d{-}2) - \zeta(2d{-}1)\bigr)
  \label{eq:master}
\end{equation}
\end{observation}

各項の評価：$2/|B_2| = 12$、$4/|B_4| = 120$、$g(132) = 5$（$137$は素数）、$4(\zeta(6) - \zeta(7)) = 0.035975$。合計：$137.035975$。

\begin{table}[h]
\centering
\caption{実験値との比較}
\begin{tabular}{@{}lc@{}}
\toprule
公式(\ref{eq:master})の値 & $137.035975$ \\
実験値~\cite{Mohr2016} & $137.035999\,084(21)$ \\
\midrule
差 & $2.4 \times 10^{-5}$ \\
相対誤差 & $1.7 \times 10^{-7}$ \\
\bottomrule
\end{tabular}
\end{table}

\subsection{構造的分解}

\textbf{第1項：1D真空（12）。} $2/|B_2| = 1/|\zeta(-1)|$。1次元カシミールエネルギーの分母。

\textbf{第2項：3D真空（120）。} $4/|B_4| = 1/\zeta(-3) = 5!$。3次元カシミールエネルギーの分母。

\textbf{第3項：素数ギャップ（5）。} $g(132) = 5$。合成数$132$から次の素数$137$までの距離。$132 = 10/|B_{10}|$であり、$2/|B_2| + 4/|B_4| = 10/|B_{10}|$はベルヌーイ数の恒等式。

\textbf{第4項：高次元補正（0.036）。} $d(\zeta(2d-2) - \zeta(2d-1))$（$d = 4$）。時空次元でスケールされた6D/7Dゼータ値差。

\section{予測と検証}

\textbf{予測1：QED走りのζ表現。} $\Delta\alpha^{-1} = (d/|\zeta(1-d-2)|)(\zeta(2d-2) - \zeta(2d-1))$。$d=4$で$\Delta\alpha^{-1} = 9.066$。実験値$9.085$。一致：$0.2\%$。

\textbf{予測2：} $\alpha^{-1}(m_Z) \approx 1/\zeta(-3) + \dim\mathrm{SU}(3) = 128$。実験値$127.951$。一致：$0.04\%$。

\textbf{予測3：} $d = 3$の仮想宇宙では$\alpha^{-1} = 137.136$（検証不能だが次元依存性を例示）。

\section{これは数秘術か？}

\subsection{探索空間の評価}
この形式の式の数は$\sim 10^4$個。ランダム式が$2 \times 10^{-5}$精度で$\alpha^{-1}$に一致する確率は$\sim 10^{-3}$（$0.1\%$）。数値一致\emph{だけ}では決定的でない。

\subsection{数秘術を超える構造的制約}
\begin{enumerate}
\item 各項に独立な物理的解釈（1D真空、3D真空、素数ギャップ、高次元補正）。
\item ベルヌーイ恒等式 $2/|B_2| + 4/|B_4| = 10/|B_{10}|$ は数学的事実。
\item 素数ギャップ $g(132) = 5$ は数論的事実。
\item 係数 $d = 4$ は観測される時空次元。フィッティングパラメータではない。
\item 公式から独立な予測（QED走り、$\alpha^{-1}(m_Z)$）が$0.2\%$と$0.04\%$で確認。
\end{enumerate}

\subsection{導出に必要なもの}
\begin{enumerate}
\item 真空エネルギーがベルヌーイ数で計算される理論的必然性（NCGスペクトル作用~\cite{CCM2007}が候補）。
\item $\lfloor\alpha^{-1}\rfloor$が素数であることを強制するメカニズム（$\Spec(\mathbb{Z})$上の安定性条件）。
\item 第4項の係数$d$がスペクトル幾何学から導かれること。
\end{enumerate}

\section{算術的時空との接続}

宇宙が$\Spec(\mathbb{Z})$上にある~\cite{paper_F}ならば：$\alpha^{-1}$は$\Spec(\mathbb{Z})$の算術幾何学（ベルヌーイ数 = ζ特殊値 = コホモロジー不変量）で決まる。素数ギャップは$\Spec(\mathbb{Z})$の離散幾何学（閉点の分布）を反映。$\alpha^{-1}$が算術で$10^{-7}$精度で表現できることは、この仮説の最も具体的な傍証。

\section{結論}

ベルヌーイ数、素数ギャップ、ゼータ値による$\alpha^{-1}$の表現を提示した。自由パラメータなしで相対誤差$1.7 \times 10^{-7}$。導出ではなく観察であるが、物理定数が$\Spec(\mathbb{Z})$の算術幾何学で決定されるという仮説と整合する構造的に豊かな結果。

\section*{謝辞}
本研究はWright Brothers共同研究グループによる独立研究として遂行された。

\begin{thebibliography}{15}
\bibitem{Mohr2016} P.~J.~Mohr et al., Rev. Mod. Phys. \textbf{88}, 035009 (2016).
\bibitem{paper_A} Wright Brothers, companion paper A (2026).
\bibitem{paper_F} Wright Brothers, companion paper F (2026).
\bibitem{CCM2007} A.~H.~Chamseddine et al., Adv. Theor. Math. Phys. \textbf{11}, 991 (2007).
\end{thebibliography}

\end{document}
